\chapter{Eye Infection}
\index{Eye Infection}

\section*{Definition and Overview}
Eye infections occur when viruses, bacteria, fungi, or other organisms infect the eye or its surrounding structures. The most common is conjunctivitis, infection or inflammation of the clear membrane over the white of the eye, often called pink eye. Other infections affect the cornea (keratitis), the eyelid (styes and blepharitis), or, rarely, the inside of the eye (endophthalmitis). Most eye infections are mild and settle on their own, but some, particularly in contact lens wearers, can threaten sight and require urgent treatment.

\section*{Epidemiology}
Conjunctivitis is one of the most common reasons for GP visits and is especially common in children, spreading quickly through child care and schools. Viral conjunctivitis is the most common type, while bacterial conjunctivitis is also frequent. Keratitis is less common but more serious and is a particular risk for contact lens wearers, who should seek care for any red, painful eye.

\section*{Aetiology and Pathophysiology}
Eye infections are caused by different organisms.
\begin{itemize}
    \item \textbf{Viral conjunctivitis:} usually adenovirus, with watery discharge, often with a cold
    \item \textbf{Bacterial conjunctivitis:} staphylococci, streptococci, and Haemophilus, with thick discharge
    \item \textbf{Keratitis:} infection of the cornea, most often from contact lens overuse, poor lens hygiene, or injury
    \item \textbf{Styes and blepharitis:} infection and inflammation of the eyelid glands
    \item \textbf{Herpes infections:} the herpes simplex virus can infect the cornea
    \item \textbf{Endophthalmitis:} infection inside the eye, which is rare and usually follows surgery or injury
    \item \textbf{Transmission:} viral and bacterial conjunctivitis spread by hand contact and shared towels
\end{itemize}

\section*{Clinical Features}
Symptoms vary with the type of infection.
\begin{itemize}
    \item Redness of the white of the eye
    \item Watery, thick, or sticky discharge
    \item Itching, burning, or a gritty feeling
    \item Crusting of the eyelids, especially on waking
    \item Pain and light sensitivity, which suggest keratitis
    \item Blurred vision, which is a warning sign
    \item In keratitis: severe pain, redness, and a corneal ulcer that may be visible
\end{itemize}

\section*{Differential Diagnosis}
Not all red eyes are infections.
\begin{itemize}
    \item Allergic conjunctivitis, with intense itching
    \item Dry eye, with a gritty feeling
    \item Foreign bodies and corneal abrasions
    \item Acute glaucoma, a serious cause of a red, painful eye
    \item Uveitis and iritis
    \item Subconjunctival haemorrhage, a harmless blood spot
\end{itemize}

\section*{When to Seek Medical Care}
See a doctor for eye pain, reduced vision, or a red eye that does not settle, and for any red or painful eye in a contact lens wearer, a newborn, or someone with a weakened immune system.

\textbf{Contact your local emergency services for:}
\begin{itemize}
    \item Sudden severe eye pain
    \item Sudden loss or change of vision
    \item A red eye with severe pain, nausea, and blurred vision
    \item A chemical injury to the eye
    \item A visible ulcer on the clear part of the eye
\end{itemize}

\section*{Diagnosis and Investigations}
Diagnosis is usually clinical.
\begin{itemize}
    \item \textbf{History and examination:} the type of discharge, symptoms, and contact lens use
    \item \textbf{Slit-lamp examination:} assesses the cornea and the depth of the problem
    \item \textbf{Fluorescein staining:} reveals corneal ulcers and abrasions
    \item \textbf{Swabs:} for severe, recurrent, or unusual infections
    \item \textbf{Referral:} to an ophthalmologist for keratitis and other serious infections
\end{itemize}

\section*{Management}
Treatment depends on the cause.
\begin{itemize}
    \item \textbf{Viral conjunctivitis:} settles on its own within one to two weeks; cool compresses and lubricating drops help
    \item \textbf{Bacterial conjunctivitis:} antibiotic drops or ointment
    \item \textbf{Keratitis:} urgent treatment with antibiotic or antiviral drops, sometimes in hospital
    \item \textbf{Styes:} warm compresses; most drain on their own
    \item \textbf{Blepharitis:} eyelid hygiene with warm compresses and cleaning
    \item \textbf{Endophthalmitis:} urgent treatment with antibiotics injected into the eye
    \item \textbf{Hygiene:} hand washing and avoiding shared towels to prevent spread
\end{itemize}

\section*{Complications}
\begin{itemize}
    \item Corneal ulcers and scarring from untreated keratitis
    \item Permanent vision loss in severe or untreated infections
    \item Spread of infection to others
    \item Recurrent infections
    \item Rarely, spread of infection within the eye
\end{itemize}

\section*{Prognosis and Outcomes}
Most eye infections resolve fully with simple treatment or on their own. Viral conjunctivitis settles within a couple of weeks, bacterial conjunctivitis responds quickly to antibiotics, and styes are minor. Keratitis and endophthalmitis are serious but treatable when urgent care is sought, which is why contact lens wearers must never ignore a red, painful eye.

\section*{Prevention and Self-Care}
\begin{itemize}
    \item Wash hands often and avoid touching the eyes
    \item Do not share towels, eye makeup, or contact lens supplies
    \item Follow contact lens hygiene: clean and replace lenses and solution as advised
    \item Remove contact lenses at the first sign of a red or painful eye
    \item Clean discharge gently with a clean cloth
    \item Do not use old eye drops
    \item Seek assessment for any eye pain or vision change
\end{itemize}

\section*{Sources and Further Reading}
The following reputable sources provide further information for healthcare students and patients seeking reliable, current advice about eye infections.
\begin{itemize}
    \item Optometry Australia. \textit{Conjunctivitis and eye infections.} https://www.optometry.org.au/
    \item Healthdirect Australia. \textit{Conjunctivitis.} https://www.healthdirect.gov.au/conjunctivitis
    \item The Royal Children's Hospital Melbourne. \textit{Conjunctivitis fact sheet.} https://www.rch.org.au/
    \item NHS. \textit{Conjunctivitis and keratitis.} https://www.nhs.uk/conditions/conjunctivitis/
    \item Mayo Clinic. \textit{Pink eye.} https://www.mayoclinic.org/diseases-conditions/pink-eye/
    \item MSD Manual. \textit{Infections of the eye.} https://www.msdmanuals.com/
\end{itemize}
